\documentclass[12pt,a4paper]{article}
\usepackage[utf8]{inputenc}
\usepackage[spanish]{babel}
\usepackage{graphicx}
\usepackage{listings}
\usepackage{xcolor}
\usepackage[margin=2.5cm]{geometry}
\usepackage{hyperref}

\definecolor{codebg}{RGB}{245,245,245}
\definecolor{codegreen}{RGB}{0,128,0}
\definecolor{codepurple}{RGB}{128,0,128}
\definecolor{codeblue}{RGB}{0,0,200}

\lstset{
    language=Java,
    backgroundcolor=\color{codebg},
    basicstyle=\ttfamily\footnotesize,
    keywordstyle=\color{codeblue}\bfseries,
    stringstyle=\color{codegreen},
    commentstyle=\color{codegreen}\itshape,
    numberstyle=\tiny\color{gray},
    numbers=left,
    numbersep=5pt,
    breaklines=true,
    frame=single,
    rulecolor=\color{gray},
    tabsize=4,
    showstringspaces=false,
    literate={á}{{\'a}}1 {é}{{\'e}}1 {í}{{\'i}}1 {ó}{{\'o}}1 {ú}{{\'u}}1 {ñ}{{\~n}}1
}

\title{\textbf{Antidoping -- Marrón 2} \\ \large Ingeniería del Software 2025--2026}
\author{Ibai Martinez}
\date{\today}

\begin{document}

\maketitle

% ============================================================
\section{Código del Antidoping}
% ============================================================

% >>> Pon aquí tu código de antidoping <<<
\textbf{Código:} \texttt{9}

% ============================================================
\section{Descripción del cambio propuesto}
% ============================================================

El cambio consiste en modificar el comportamiento del botón inferior (\textit{Book Flight}) de la interfaz gráfica para que, al pulsarlo, muestre en el desplegable (\texttt{JComboBox}) \textbf{todos los vuelos concretos que tienen como destino la ciudad de llegada (Arrival City) seleccionada}, independientemente de:

\begin{itemize}
    \item La \textbf{fecha} seleccionada.
    \item La \textbf{ciudad de salida} (Departing City) seleccionada.
\end{itemize}

Además, el botón debe estar \textbf{siempre habilitado}, sin necesidad de haber pulsado previamente ``Look for Concrete Flights'' ni de haber seleccionado un vuelo o tipo de asiento.

% ============================================================
\section{Implementación}
% ============================================================

Los cambios se realizaron en dos ficheros:

\subsection{FlightManager.java (Capa de lógica de negocio)}

Se añadió un nuevo método \texttt{getAllConcreteFlightsTo(String arrivingCity)} que realiza una consulta JPA directa a la base de datos ObjectDB para obtener todos los \texttt{ConcreteFlight} cuyo vuelo tiene como ciudad de llegada la indicada:

\begin{lstlisting}
public Collection<ConcreteFlight> getAllConcreteFlightsTo(String arrivingCity) {
    TypedQuery<ConcreteFlight> query = db.createQuery(
        "SELECT cf FROM ConcreteFlight cf WHERE cf.flight.arrivingCity = :city",
        ConcreteFlight.class);
    query.setParameter("city", arrivingCity);
    return query.getResultList();
}
\end{lstlisting}

Este método navega por la relación \texttt{cf.flight.arrivingCity} en JPQL, lo que permite obtener todos los vuelos concretos a esa ciudad sin importar la fecha ni la ciudad de origen.

\subsection{FlightBooking.java (Capa de presentación)}

Se reescribió la lógica del botón \texttt{bookFlight}:

\begin{lstlisting}
bookFlight = new JButton("Book Flight");
bookFlight.addActionListener(new ActionListener() {
    public void actionPerformed(ActionEvent e) {
        // Mostrar todos los vuelos a la ciudad de llegada seleccionada
        flightInfo.removeAllElements();
        fareButtonGroup.clearSelection();
        bussinesTicket.setEnabled(false);
        firstTicket.setEnabled(false);
        touristTicket.setEnabled(false);
        selectedConcreteFlight = null;

        Object arrive = arrivalCity.getSelectedItem();
        if (arrive != null) {
            Collection<ConcreteFlight> allFlights =
                businessLogic.getAllConcreteFlightsTo(arrive.toString());
            for (ConcreteFlight cf : allFlights) {
                flightInfo.addElement(cf);
            }
            if (allFlights.isEmpty()) {
                searchResult.setText("No flights to " + arrive.toString());
            } else {
                searchResult.setText("All flights to " + arrive.toString() + ":");
            }
        } else {
            searchResult.setText("Select an arrival city.");
        }
    }
});
\end{lstlisting}

Cambios clave:
\begin{itemize}
    \item Se \textbf{eliminó toda la lógica anterior} de reserva (\texttt{bookConcreteFlight}) del botón.
    \item Se \textbf{eliminaron todas las llamadas a} \texttt{bookFlight.setEnabled(false)} para que el botón esté siempre activo.
    \item Al pulsar, se lee la \textit{Arrival City} del combo y se llama a \texttt{getAllConcreteFlightsTo()}, que devuelve todos los vuelos concretos a esa ciudad.
    \item Los resultados se cargan en el \texttt{JComboBox} (\texttt{flightInfo}).
\end{itemize}

\subsection{Datos de prueba modificados}

Para verificar que el cambio funciona con \textbf{distintas fechas}, se modificaron las fechas de algunos vuelos concretos de los vuelos F2 y F5 (ambos con destino Madrid):

\begin{lstlisting}
// F2: Donostia -> Madrid (fechas variadas)
f2.addConcreteFlight("CF2-1", newDate(2023,2,15), 1,2,3, "12:00"); // Marzo
f2.addConcreteFlight("CF2-2", newDate(2023,3,10), 3,0,3, "12:30"); // Abril
f2.addConcreteFlight("CF2-3", newDate(2023,4,5),  1,2,2, "13:00"); // Mayo
f2.addConcreteFlight("CF2-6", newDate(2023,5,20), 3,3,0, "16:00"); // Junio
f2.addConcreteFlight("CF2-7", newDate(2023,6,1),  3,3,0, "17:00"); // Julio

// F5: Sevilla -> Madrid (fechas variadas)
f5.addConcreteFlight("CF5-1", newDate(2023,2,22), 1,2,3, "8:00");  // Marzo
f5.addConcreteFlight("CF5-4", newDate(2023,4,23), 3,3,0, "14:00"); // Mayo
\end{lstlisting}

% ============================================================
\section{Capturas de pantalla}
% ============================================================

\begin{figure}[h!]
    \centering
    \includegraphics[width=0.8\textwidth]{1.png}
    \caption{Interfaz inicial: el botón ``Book Flight'' está habilitado por defecto, sin necesidad de haber buscado vuelos previamente.}
\end{figure}

\begin{figure}[h!]
    \centering
    \includegraphics[width=0.8\textwidth]{2.png}
    \caption{Tras pulsar ``Book Flight'': el desplegable muestra todos los vuelos concretos a la ciudad de llegada seleccionada, con distintas fechas y distintas ciudades de origen.}
\end{figure}

\end{document}
